\section{Class Declarations}
When coding \Java classes, interfaces, enums, records, and annotations, you should follow the guidelines below:


\subsubsection{Principles}
\begin{enumerate}[label=P\arabic*.]
\item{.}
\item{.}
\item{.}
\end{enumerate}








\begin{itemize}
\item{No space between a method or constructor name and the opening parenthesis “(” starting its formal parameters list. Basically the parenthesis in the declaration are “paramater parenthesis” and should be treated along the same lines as for method and constructors calls; refer to chapter “\nameref{sec:ParameterParenthesis}” on page~\pageref{sec:ParameterParenthesis}.

Example:
\begin{lstlisting}
// WRONG!!
public abstract void myMethod (final int param);

// CORRECT
public abstract void myMethod( final int param );
\end{lstlisting}}

\item{If using generics, no space between the name of the declared class or interface and the opening “\textless”, and no blanks between the angle brackets, only around the keywords \lstinline|super| and \lstinline|extends|. Same when using the class or interface later.

Example:
\begin{lstlisting}
// WRONG
public interface MyInterface < T >
{ 
    … 
}
//  interface MyInterface

public OtherInterface < T, R > extends Function < T, R >, MyInterface < T >
{
    …
    
    public < U super T > U myMethod( final T arg );
}
//  interface OtherInterface

public MyClass implements otherInterface < Instant, String >
{
    …
}
//  class MyClass

// CORRECT
public interface MyInterface<T>
{ 
    … 
}
//  interface MyInterface

public OtherInterface<T,R> extends Function<T,R>, MyInterface<T>
{
    …
    
    public <U super T> U myMethod( final T arg );
}
//  interface OtherInterface

public MyClass implements otherInterface<Instant,String>
{
    …
}
//  class MyClass
\end{lstlisting}}

\item{An opening curly brace “\{” appears at a line by itself on the same indentation level as the declaration statement; anything afterwards will be indented with one additional level.

The corresponding closing curly brace “\}” starts a new line, indented to match its corresponding opening statement. For classes or interfaces, the name of the class or interface is repeated in a new line after the closing brace.

Example:\footnote{I omitted the structuring comments here; refer to chapters \tqfullvref{sec:ClassAndInterfaceDeclarations} and \tqfullvref{sec:StructuringComments}. And the bad sample is \textit{really} bad!}
\index{java.time!Instant!now()}\index{java.time!Instant!toString()}
\begin{lstlisting}
// AVOID!!
public final class MyClass {
public static InnerClass {
public final String innerMethod {
return Instant.now().toString(); } }
public static final void main( final String... args ) {
System.out.println( new InnerClass().innerMethod() ); } }

// RECOMMENDED
public final class MyClass 
{
    public static InnerClass 
    {
        public final String innerMethod 
        {
            return Instant.now().toString(); 
        }   //  innerMethod()
    }
    //  class InnerClass
    
    public static final void main( final String... args ) 
    {
        System.out.println( new InnerClass().innerMethod() ); 
    }   // main()
}
//  class MyClass    
\end{lstlisting}
Both versions compile properly, can be executed and will return the same result ~…}

\item{Inner classes, constants, attributes, constructors, methods, they all are separated from each other by a  single blank line (see chapter \tqfullref{sec:BlankLines}.}
\end{itemize}


%==============================================================================
% End of file!!
%==============================================================================
